% Answers to Chapter 9, from lectures/L09/appendix/c_solutions.md; the self-test from
% lectures/L09/README.md.

\facitavsnitt{c9}{Kapitel 9}{Trigonometriska funktioner}
\begin{facit}

\svar{1.1} \sv{a} $0\,\text{V}$
\sv{b} $0 \leq u_{\text{c}}(t) \leq 5(1 - e^{-5}) \approx 4{,}97\,\text{V}$
\sv{c} $t = -2\ln 0{,}2 \approx 3{,}22\,\text{s}$

\svar{1.2} \sv{a} $0^{\circ}$, $6^{\circ}$, $8^{\circ}$, $6^{\circ}$ och $0^{\circ}$
\sv{b} En nedåtvänd parabel genom punkterna i a), symmetrisk kring $t = 4\,\text{s}$; den största
vinkeln, $8^{\circ}$, nås vid $t = 4\,\text{s}$.
\sv{c} $t = 1\,\text{s}$ och $t = 7\,\text{s}$

\svar{1.3} \sv{a} $x^2 + x - 12 = (x + 4)(x - 3)$, $x^2 - 16 = (x + 4)(x - 4)$
\sv{b} $x = 4$ och $x = -4$
\sv{c} $G(x) = \frac{x - 3}{x - 4}$, $x \neq \pm 4$

\svar{2.1} \sv{a} $-\frac{\pi}{6}$ \sv{b} $\frac{\pi}{9}$ \sv{c} $\frac{5\pi}{4}$
\sv{d} $-\frac{\pi}{4}$ \sv{e} $\frac{5\pi}{3}$

\svar{2.2} \sv{a} $90^{\circ}$ \sv{b} $135^{\circ}$ \sv{c} $-150^{\circ}$
\sv{d} $\approx -57{,}3^{\circ}$ \sv{e} $\approx 126{,}1^{\circ}$

\svar{2.3} $\abs{U} = 8\,\text{V}$;\enspace $f = 10\,\text{Hz}$;\enspace
$\omega = 20\pi \approx 62{,}8\,\text{rad/s}$;\enspace $\delta = \frac{\pi}{3}\,\text{rad}$;\enspace
$u(t) = 8\sin(20\pi t + \frac{\pi}{3})\,\text{V}$

\svar{2.4} $u(t) = 3\sin(50\pi t + \frac{\pi}{4})\,\text{V}$. Kurvan har amplituden $3\,\text{V}$
och periodtiden $40\,\text{ms}$ och ligger $T/8 = 5\,\text{ms}$ före en ofasad sinuskurva: toppen
nås vid $t = 5\,\text{ms}$.

\svar{2.5} $\delta \approx -1{,}63\,\text{rad}$ eller $\delta \approx 1{,}00\,\text{rad}$

\svar{3.1} \sv{a} $\frac{\pi}{3}$ \sv{b} $\frac{2\pi}{3}$ \sv{c} $-\frac{\pi}{2}$
\sv{d} $\frac{7\pi}{6}$ \sv{e} $\frac{\pi}{12}$ \sv{f} $2\pi$

\svar{3.2} \sv{a} $60^{\circ}$ \sv{b} $225^{\circ}$ \sv{c} $-90^{\circ}$ \sv{d} $210^{\circ}$
\sv{e} $\approx 171{,}9^{\circ}$ \sv{f} $\approx 28{,}6^{\circ}$

\svar{3.3} \sv{a} $0$ \sv{b} $1$ \sv{c} $0$ \sv{d} $-1$ \sv{e} $\frac{1}{2}$
\sv{f} $\frac{\sqrt{2}}{2} \approx 0{,}707$

\svar{3.4} Uppifrån: $T = 20\,\text{ms}$, $\omega = 100\pi \approx 314\,\text{rad/s}$;\enspace
$f = 1\,\text{kHz}$, $\omega = 2000\pi \approx 6283\,\text{rad/s}$;\enspace
$f = 200\,\text{Hz}$, $T = 5\,\text{ms}$;\enspace
$T = 2{,}5\,\text{ms}$, $\omega = 800\pi \approx 2513\,\text{rad/s}$

\svar{3.5} I ordningen $\abs{U}$, $\omega$, $f$, $T$, $\delta$:
\sv{a} $5\,\text{V}$, $100\pi\,\text{rad/s}$, $50\,\text{Hz}$, $20\,\text{ms}$, $0$
\sv{b} $12\,\text{V}$, $200\pi\,\text{rad/s}$, $100\,\text{Hz}$, $10\,\text{ms}$, $\frac{\pi}{2}$
\sv{c} $3\,\text{V}$, $50\pi\,\text{rad/s}$, $25\,\text{Hz}$, $40\,\text{ms}$, $-\frac{\pi}{6}$
\sv{d} $8\,\text{V}$, $2000\pi\,\text{rad/s}$, $1000\,\text{Hz}$, $1\,\text{ms}$, $0$

\svar{3.6} \sv{a} $u(t) = 10\sin(100\pi t)\,\text{V}$
\sv{b} $u(t) = 6\sin(200\pi t + \frac{\pi}{2})\,\text{V}$
\sv{c} $u(t) = 2\sin(400\pi t - \frac{\pi}{3})\,\text{V}$
\sv{d} $u(t) = 325\sin(100\pi t + \frac{\pi}{4})\,\text{V}$

\svar{3.7} \sv{a} $0\,\text{V}$ \sv{b} $10\,\text{V}$ \sv{c} $0\,\text{V}$ \sv{d} $-10\,\text{V}$
\sv{e} $t = 0$, $10$ och $20\,\text{ms}$

\svar{3.8} \sv{a} $\approx 8{,}49\,\text{V}$ \sv{b} $u(t) = 325\sin(100\pi t)\,\text{V}$
\sv{c} $0{,}72\,\text{W}$

\svar{3.9} \sv{a} $5\,\text{ms}$ \sv{b} $\approx 1{,}67\,\text{ms}$
\sv{c} $\delta = \frac{\pi}{5} \approx 0{,}628\,\text{rad}$
\sv{d} $\delta = -\frac{\pi}{2}\,\text{rad}$

\svar{3.10} \sv{a} $u_2$ ligger före $u_1$ med $60^{\circ}$. \sv{b} $\approx 3{,}33\,\text{ms}$
\sv{c} $u_1 = 4\,\text{V}$, $u_2 = 2\,\text{V}$
\sv{d} Ja: båda har amplituden $4\,\text{V}$ och frekvensen $50\,\text{Hz}$; bara fasen skiljer.

\svar{3.11} \sv{a} $t = \frac{1}{600}\,\text{s} \approx 1{,}67\,\text{ms}$
\sv{b} $t = \frac{5}{600}\,\text{s} \approx 8{,}33\,\text{ms}$ \sv{c} $t = 5\,\text{ms}$
\sv{d} Sinus är periodisk: varje lösning upprepas en gång per period, alltså var
$20$:e millisekund.

\svar{3.12} \sv{a} Falskt: värdemängden är $[-1, 1]$.
\sv{b} Sant: $T = \frac{1}{f}$.
\sv{c} Falskt: $\omega$ mäts i rad/s; det är frekvensen $f$ som mäts i hertz.
\sv{d} Sant: kurvan når sina värden tidigare.
\sv{e} Sant: ett halvt varv är $\pi$ radianer.
\sv{f} Falskt: amplituden mäts från noll till toppvärdet; från lägsta till högsta värde är
$2\abs{U}$.

\svar[Testa dig själv]{T} \sv{1} $150^{\circ} = \frac{5\pi}{6}$ och $\frac{5\pi}{6} = 150^{\circ}$
\sv{2} $u(t) = 5\sin(100\pi t + \frac{\pi}{4})\,\text{V}$

\end{facit}
